\documentclass{article}[12]

% CREDIT TO Ioannis Spyridopoulos, Nick Frazier, Margaret McKeehan,
%           and Mark Agerton at Rice University for template.
%           Also, to Jason Blevins <jrblevin@sdf.lonestar.org> for 
%           initial template at 
%               http://jblevins.org/projects/cv-template/
% THANKS TO Jacquie Frazier for helpful suggestions



% Preamble has all of the relevant formatting info
%!TEX root = ./CV Template 2.2.tex

% ---------- Packages ------------

\usepackage{geometry}                % for margins
%\usepackage{tgpagella}               % font 
\usepackage{titlesec}                % format section titles
\usepackage{pbox}                    % boxes for diss. committee
\usepackage{fancyhdr}                % header + footer
\usepackage[yyyymmdd]{datetime}      % for date format at end
\usepackage{tabularx}                % to get tables to fill width
% \usepackage{showframe}               % uncomment to see margins
\usepackage{xcolor}                  % to define the Rice Univ blue
\usepackage{enumitem}                % Customize itemize environments
\usepackage{soul}                    % Allows underlining

\usepackage{hyperref}                % links for email + web

% ---------------------------------------------------------
% ----------             FORMATTING            ------------
% ---------------------------------------------------------

% format of section heading
% see http://www.ctex.org/documents/packages/layout/titlesec.pdf
\titleformat{\section}                       % command
    [display]                                % shape
    {\Large\scshape}                         % format
    {}                                       % label
    {0ex}                                    % sep
    {\vspace{-4.0ex plus .2ex minus .2ex}}   % code preceding the title body
    [\vspace{-2.8ex}\rule{\textwidth}{0.3pt}\vspace{-1.5ex}]  % code following the title body

% margins
\geometry{
  body={6.5in, 9.0in},                       % adjust right + bottom here
  left=1.0in,
  top=1.0in
}

% formatting for all itemize
\setlist[itemize]{ 
        leftmargin  = 0.0em, 
        itemsep     = 0.0em, 
        topsep      = 0.0em,
        label       = {},
        parsep      = 0.5em
}

% formatting for second level of itemize
\setlist[itemize,2]{
        leftmargin  = 2.0em, 
        itemsep     = 0.0em
}

% -------------- format header & footer ------------

% get header + footer 
\pagestyle{fancy}

% suppress line drawn by default by fancyhdr
\renewcommand{\headrulewidth}{0pt}

% shows date compiled
% \cfoot{(Updated: \mydate\today)}
\cfoot{}

% your name on the right
\rhead{\emph{\myshortname, \  p. \thepage}}
\lhead{}

\fancypagestyle{laststyle}{
   \fancyfoot[C]{(Updated \mydate\today)}
}

% ---------- Margins and other formatting ------------

% space b/w paragraphs
\setlength{\parskip}{1.1ex}

% default to no indents
\setlength\parindent{0pt}

% format for date at the end
\newdateformat{mydate}{\monthname[\THEMONTH] \THEDAY, \THEYEAR}

% The Rice Univ Blue
\definecolor{RiceBlue}{RGB}{0,82,147}

% so that URLs are not monospace
\urlstyle{same}

% macro for email
\newcommand{\email}[1]{\href{mailto:#1}{#1}}

% macro to force indent of a line
\newcommand{\forceindent}{\leavevmode{\parindent=24pt\indent}}




% ---------------------------------------------------------
% -----------         YOUR INFORMATION          -----------
% ---------------------------------------------------------

\newcommand{\myname}{Zhiran Qin}
\newcommand{\myshortname}{Z. Qin}
\newcommand{\myemail}{zhrqin@ucdavis.edu}
\newcommand{\mywebsite}{https://cliobserver.github.io/zhiran-AgEcon.github.io/}
\newcommand{\mywebsiteNoHttp}{https://cliobserver.github.io/zhiran-AgEcon.github.io/}
\newcommand{\myphone}{(530) 564-9235}
\newcommand{\myaddress}{
Dept. of Agricultural and Resource Economics PhD-27 \\ 
University of California, Davis \\ 
Davis, CA 95616}
\newcommand{\mykeywords}{economics, econometrics, macroeconomics}

% ---------------------------------------------------------
% -----------    URL COLOR + PDF METADATA       -----------
% ---------------------------------------------------------

% URL and link colors
\newcommand{\linkcolor}{RiceBlue}     % can change to black

\hypersetup{
            colorlinks  = true, 
            linkcolor   = {\linkcolor},
            filecolor   = {\linkcolor},
            urlcolor    = {\linkcolor},
            pdfauthor   = {\myname},
            pdfkeywords = {\mykeywords},
            pdfsubject  = {Curriculum Vitae},
            pdftitle    = {\myname: Curriculum Vitae},
            pdfpagemode = UseNone,
}


% ---------------------------------------------------------
% ----------              RESUME               ------------
% ---------------------------------------------------------

\begin{document}

% suppress header + footer on p. 1
\thispagestyle{empty}

% Name
\begin{center}
{\huge\scshape\myname}
\end{center}

\vspace{3em}

% Contact info
\begin{minipage}[t]{0.48\textwidth}
    \href{\mywebsite}{\mywebsiteNoHttp} \\
    \email{\myemail}  \\ 
    \myphone 
\end{minipage}
\hfill
\begin{minipage}[t]{0.48\textwidth}
    \begin{flushright}
    \myaddress
    \end{flushright}
\end{minipage}


% ----------------------------------------
\section{Education}
% ----------------------------------------

\begin{itemize}[labelindent=-25pt,itemindent=-20pt,leftmargin=!]
    \item Ph.D. in Agricultural and Resource Economics, UC Davis  \hfill (Anticipated) June 2027  \\
    \textsc{Dissertation}: \emph{Three Essays in Agricultural Economics}

    %\item M.A. in Economics, Rice University  \hfill (Anticipated) May 2017 

    \item B.A. in Agricultural Economy and Management, Renmin University of China \hfill June 2021 
\end{itemize}

% % ----------------------------------------
% \section{Dissertation Committee (pbox 1 row)}
% % ----------------------------------------

% \pbox{0.33\textwidth}{
% Prof. Pierre Mérel (\textbf{Co-chair}) \\
% Department of Agricultural and Resource Economics \\
% University of California, Davis \\
% (713) 348-1111\\
% \href{mailto:arrow@rice.edu}{arrow@rice.edu}
% } %
% %
% \hfill
% %
% \pbox{0.33\textwidth}{
% Prof. Friedrich Hayek \\
% Department of Economics \\
% Rice University \\
% (713) 348-2222\\
% \href{mailto:hayek@rice.edu}{hayek@rice.edu}
% }%
% %
% \hfill
% %
% %
% \pbox{0.33\textwidth}{
% Prof. John Hicks \\
% Department of Economics\\
% Rice University \\
% (713) 348-3333 \\
% \href{mailto:hicks@rice.edu}{hicks@rice.edu}
% }

% ----------------------------------------
\section{Dissertation Committee}
% ----------------------------------------


\begin{minipage}{.5\textwidth}
Prof. Pierre M\'erel (\textbf{Co-chair}) \\
Department of Agricultural and Resource Economics \\
University of California, Davis\\
%(713) 348-1111\\
\href{mailto:prmerel@ucdavis.edu}{prmerel@ucdavis.edu}
\end{minipage}%
%
%
\begin{minipage}{0.5\textwidth}
Distinguished Prof. Richard J. Sexton (\textbf{Co-chair})\\
Department of Agricultural and Resource Economics \\
University of California, Davis \\
%(713) 348-2222\\
\href{mailto:rjsexton@ucdavis.edu}{rjsexton@ucdavis.edu}
\end{minipage}

% CODE TO INSERT A \vspace EQUIVALENT of 
%    "The normal vertical distance between lines in a paragraph."
\vspace{1\baselineskip plus 10pt minus 10pt} 

\begin{minipage}{0.5\textwidth}
Assistant Prof. James E. Sayre \\
Department of Agricultural and Resource Economics \\
University of California, Davis \\
%(713) 348-3333 \\
\href{mailto:jsayre@ucdavis.edu }{jsayre@ucdavis.edu }
\end{minipage}

% ----------------------------------------
\section{Research Fields}
% ----------------------------------------

\begin{itemize}
    \item Primary: Agricultural and Trade Econometrics
    \item Secondary: Industrial Organization
\end{itemize}


% ----------------------------------------
\section{Academic Publications}
% ----------------------------------------

\begin{itemize}
    \item ``Policy-induced expansion of organic farmland: implications for food prices and welfare.'' (2023) \url{https://doi.org/10.1093/erae/jbad024} \emph{European Review of Agricultural Economics}, 50(4), 1583--1631. (with Pierre M\'erel and Richard J. Sexton)
    \item ``A `Sideways' Supply Response in California Winegrapes.'' (2022) \emph{Journal of Wine Economics}, 17(1), 42--63. \url{https://doi.org/10.1017/jwe.2021.26} (with Sarah Consoli, Elizabeth A. Fraysse, Natalya Slipchenko, Yi Wang, Jahon Amirebrahimi, Neil Yazma, and Travis J. Lybbert)
\end{itemize}

% ----------------------------------------
\section{Research Papers}
% ----------------------------------------

\begin{itemize}[itemsep=1.5em]

\item \emph{Nutrition Transition and Population Growth: Implications for}
 \hfill \textbf{Job Market Paper} \\ \emph{Land Use and Food Prices}

The nutrition transition happening in the development world is putting pressure on world agricultural production. As countries become wealthier, people are adopting a more diversified diet, shifting from dominantly starch staples to vegetables, fruits, and animal protein. The goal of this paper is to study the implications of the nutrition transition driven by income growth in developing countries. Drawing inspiration from a new strand of modeling approaches proposed by Costinot, Donaldson, and Smith (2016), we build up a succinct, transparent, yet powerful partial equilibrium (PE) model capable of analyzing isolated or combined effects of nutrition transition and population growth under different counterfactual cases. The model yields country-specific equilibrium projections, eg, to 2050, including food and crop prices and land use changes, for each supply-demand scenario, offering clear insights into impacts of possible scenarios and policy responses. 
%\href{http://my-great-site.com/files/somepaper.pdf}{Download}

\item \emph{From Trials to Commercialization: Global Market Effects of China’s Genetically Modified Soybean and Corn Policies}
%\\ (\emph{Revise and Resubmit, \emph{Journal of Latin Cliches}}) with J.M. Keynes, LSE
\\ (\emph{Working Paper}) with Zhiyao Ma and Pierre M\'erel

China's transition from prudent pilots to nationwide commercialization of genetically modified (GM) corn and soybeans represents a potentially consequential shift in global crop supply, given the country's outsized role in world feedgrain and oilseed demand and trade. How a domestic GM rollout transmits through international markets to change production, trade flows, and welfare outcomes is a central but under-studied question, with particular relevance for export-dependent producers and import-dependent consumers. Two foundations motivate the analysis. First, China is actively scaling GM corn and soybean cultivation under a national food-security strategy that prioritizes ``bio-breeding''. Second, China's large demand profile and bilateral trade shares imply that any supply-side shock originating in its domestic agriculture transmits to world prices and welfare in proportion to that exposure. We develop and calibrate a partial equilibrium model of world crop demand and supply, disciplined by observable production, consumption, and trade shares together with behavioral elasticities. We apply a 20\% yield improvement to China's soybean acreage and obtain three main findings. The domestic supply response is large in percentage terms but transmits weakly to consumer-relevant prices, because imports supply the majority of baseline Chinese soybean volume. World-price spillovers are modest, with the largest non-China producer-price decline below 1\%. Welfare gains are concentrated in China; Brazil and the United States bear the largest producer losses, though domestic consumer gains partially offset these at the national welfare level. 
%\href{http://my-great-site.com/files/somepaper.pdf}{Available on SSRN}

\item \emph{Contract Prices for Agricultural Products When Buyers Have Heterogeneous Valuations and Binding Capacity Constraints}
%\\ (\emph{Revise and Resubmit, \emph{Journal of Latin Cliches}}) with J.M. Keynes, LSE
\\ (\emph{Working Paper}) with Richard J. Sexton

This paper examines contract price determination in agricultural product procurement markets characterized by buyer concentration, heterogeneous buyer valuations, and binding capacity constraints. Using illustrative duopsony examples and a Hotelling model framework, we show that binding capacity constraints severely attenuate competitive bidding even when multiple non-collusive buyers are present. The key mechanism is that capacity-constrained buyers internalize the profit they can guarantee on high-valuation lots, which limits their willingness to bid on lots where they are not the high-valuation buyer. This, in turn, reduces the competitive pressure faced by high-valuation buyers, pushing equilibrium contract prices toward sellers' outside options rather than buyers' marginal value products. The result holds regardless of whether high-valuation lots are sufficient to meet each buyer's capacity. When total production exceeds total buyer capacity, each packer effectively behaves as a local monopsony, retaining the bulk of market surplus. Excess capacity reverses this outcome by restoring competitive pressure. The analysis highlights capacity constraints as a previously underappreciated mechanism through which market power is exercised in agricultural procurement markets, with implications for farm prices, farmer welfare, and antitrust policy. 


\end{itemize}

% ----------------------------------------
\section{Scholarships, Honors, and Awards}
% ----------------------------------------

\begin{itemize}   
    \item JASTRO-SHIELDS Graduate Awards \hfill 2024, 2025
\end{itemize}

% ----------------------------------------
% \vspace*{-1.2cm}  % put above section title if at top of pg
\section{Teaching Fields}
% ----------------------------------------

Microeconomics, Econometrics


% ----------------------------------------
\section{Teaching Experience}
% ----------------------------------------

\begin{itemize}[itemsep=0.5em]
    
    \item \textbf{Teaching Assistant}, UC Davis
    \begin{itemize}
        \item Graduate: Microeconomic Analysis (204A/B); Applied Econometrics (256B) %\hfill Fall 2023-Present 
        \item Undergraduate: Futures \& Options Markets (139); Operations Research \& Management Science (155); Principles of Finance (171) %\hfill Fall 2010-Fall 2012
    \end{itemize}
    
    \item \textbf{Instructor}, UC Davis
    \begin{itemize}
        \item Undergraduate: Intermediate Microeconomics: Imperfect Competition, Markets \& Welfare Economics (100B)  \hfill Summer 2024
    \end{itemize}

\end{itemize}


% ----------------------------------------
\section{Professional Activities}
% ----------------------------------------

\begin{itemize}[itemsep=0.75em]

        \item \textbf{Conferences}
    \begin{itemize}
        %\item ``Sports in the Age Sports," 2023 AAEA, Washington DC \hfill 2023
        \item ``Nutrition Transition and Population Growth: Implications for Land Use and Food Prices," 2025 AAEA, Denver, CO \hfill 2025
        \item ``From Trials to Commercialization: Global Market Effects of China’s Genetically Modified Soybean and Corn Policies," 2026 AAEA, Kansas City, MO \hfill 2026
        \item ``Contract Prices for Agricultural Products When Buyers Have Heterogeneous Valuations and Binding Capacity Constraints," 2026 AAEA, Kansas City, MO \hfill 2026
    \end{itemize}

%    \item \textbf{Referee for} Journal of Obligatory Economics, Review of Useless Economics
    
    % \item \textbf{Association of Useless Facts}
    % \begin{itemize}
    %     \item President, Rice Univ Chapter \hfill 2010--2012
    %     \item Student representative, National chapter \hfill 2013
    % \end{itemize}
    
    % \item \textbf{Invited presentations}
    % \begin{itemize}
    %     \item Luna University, The moon, ``Economics of space travel'' \hfill Sept 2013
    %     \item Extraterrestris institute, Alpha Centauri, ``Are we really all dead in the long run?'' \hfill Jan 2101
    % \end{itemize}

\end{itemize}

% % ----------------------------------------
% \section{Other Writing and Media Appearances}
% % ----------------------------------------

% \begin{itemize}[itemsep=0.5em]
%     \item Quoted on Planet Money, ``\href{https://www.marketplace.org/2017/05/26}{Will work for money: A conversation with David Ricardo}'' (May 26, 2017)

%     \item Rice University Center for Energy Studies Blog on Forbes.com
%         \begin{itemize}
%             \item \href{https://www.forbes.com/sites/thebakersinstitute/2017/05/10/the-cycle-not-automation-is-keeping-oil-gas-hiring-down}{``The Cycle, Not Automation, Is Keeping Oil \& Gas Hiring Down''} (May 10, 2017)
%             \item \href{https://www.forbes.com/sites/thebakersinstitute/2015/04/27/oil-price-exchange-rates-and-the-convoluted-impact-of-sanctions-on-russia}{``Oil Price, Exchange Rates, and the Convoluted Impact of Sanctions on Russia''} (April 27, 2015)
%         \end{itemize}

% \end{itemize}

% ----------------------------------------
\section{Other}
% ----------------------------------------

\begin{itemize}[itemsep=0.5em]
    %\item Fluent in Greek
    \item Software: R, Stata, MATLAB, Python, ArcGIS, \LaTeX
    \item Hobbies: skiing, hiking, gardening, cooking
    %\item Prior professional experience: Central Banking
    %\item Green card holder
\end{itemize}


% Ensures "updated" date comes at the end
\thispagestyle{laststyle}

\end{document}